\section{Retrieval error on the admissible domain}
\label{sec:constrained}

Proposition~1 and the transfer theorem of Section~\ref{sec:theory} describe an entry at
which the transmitted flux is positive and $q = 1 - \rho Y_4$ is bounded away from zero.
The all-band statistics of Table~\ref{tab:tails} include entries at which neither holds,
and there the inverse is not merely ill-conditioned but undefined. This section separates
the two, on a domain fixed before any model is scored.

Admissibility is a property of the tabulation. Over all $23{,}313\times285 = 6{,}644{,}205$
tabulated entries, the transmitted flux $t = Y_2 + Y_3$ is non-positive at 438 entries, with
a minimum of $-4.5\times10^{-20}$, and falls below $10^{-12}$ at 32{,}130. The spherical
albedo lies outside $[0,1)$ at 10{,}312 entries: 9{,}858 are negative, with a minimum of
$-0.0724$, and 454 are at or above one, of which 214 are exactly $2$, a fill marker rather
than an albedo. At 448 entries $q \le 0$, with a minimum of $-0.4$. Writing
$\mathcal{A} = \{t > 0,\ 0 \le Y_4 < 1\}$, the inadmissible set has 10{,}356 entries
($0.156\%$); it occupies 227 of the 285 bands and concentrates in the water-vapour band
between 1841 and 1908\,nm and at the long-wave edge between 2485 and 2493\,nm. Of the
23{,}313 states, 18{,}859 are admissible at every band.

\IfFileExists{table_domain_audit.tex}{%
\begin{table}[t]
\centering
\begin{tabular}{lrr}
\toprule
condition on the tabulated truth & entries & share \\
\midrule
transmitted flux $t=Y_2+Y_3 \le 0$ & 438 & $0.0066\%$ \\
$t < 10^{-12}$ & 32,130 & $0.4836\%$ \\
albedo $Y_4 < 0$ & 9,858 & $0.1484\%$ \\
albedo $Y_4 \ge 1$ & 454 & $0.0068\%$ \\
albedo exactly $2$ (fill marker) & 214 & $0.0032\%$ \\
\midrule
admissible $\mathcal{A}=\{t>0,\ 0\le Y_4<1\}$ & 6,633,849 & $99.8441\%$ \\
\bottomrule
\end{tabular}

\caption{Entries of the tabulated truth at which the retrieval hypotheses fail, counted
over all $6{,}644{,}205$ sample--band entries before any model is fitted.}
\label{tab:domain}
\end{table}}{}

Admissibility alone does not make the inverse stable, because $t$ can be positive and far
below the scale on which $L - Y_1$ is resolved. We therefore fix a conditioning floor in
advance, $\tau = 10^{-12}$ on $t$ and $\kappa = 0.3$ on $q$, and report the retained
coverage beside every number computed on it. On the exact tabulated components the three
regimes separate:

\begin{center}
\begin{tabular}{lrr}
\toprule
round trip at exact components & reflectance RMSE & coverage \\
\midrule
all bands & $7.38\times10^{-3}$ & $100\%$ \\
admissible, $t>0$ and $0\le Y_4<1$ & $4.27\times10^{-3}$ & $99.844\%$ \\
with the conditioning floor $t\ge\tau$, $q\ge\kappa$ & $1.08\times10^{-9}$ & $99.516\%$ \\
\bottomrule
\end{tabular}
\end{center}

The reflectance error of the exact-component round trip is therefore not a property of the
inversion on its domain of definition. It is carried by the $0.48\%$ of entries below the
floor, where the reflectance is unidentifiable or nearly so and a returned inverse measures
cancellation rather than emulator accuracy. A reflectance statistic taken over all bands is
dominated in its tail by that set, and a comparison between emulators conducted over all
bands compares, in its tail, their behaviour where no emulator can be correct.

\IfFileExists{table_constrained_retrieval.tex}{%
\begin{table}[t]
\centering
\resizebox{\textwidth}{!}{\begin{tabular}{lrrrrrr}
\toprule
 & radiance & \multicolumn{2}{c}{all bands} & \multicolumn{2}{c}{admissible domain} \\
\cmidrule(lr){3-4}\cmidrule(lr){5-6}
model & rel.\ $L^2$ [\%] & median [pp] & 95th [pp] & median [pp] & 95th [pp] \\
\midrule
cubic ridge & 0.7543 & 0.480 & 70.68 & 0.473 & 70.629 \\
network $3\times512$ & 0.1741 & 0.111 & 37.50 & 0.111 & 24.109 \\
network $+$ residual kernel & 0.0938 & 0.065 & 25.54 & 0.065 & 15.650 \\
input-scaled Mat\'ern kernel & 0.0853 & 0.028 & 24.38 & 0.027 & 16.830 \\
kernel on learned features & 0.0885 & 0.037 & 4.19 & 0.037 & 3.117 \\
convex stack & 0.0709 & 0.046 & 25.67 & 0.046 & 15.549 \\
\bottomrule
\end{tabular}
}
\caption{Absolute reflectance error at $\rho = 0.7$ on the test block of one split, over
all bands and on the admissible domain with the conditioning floor $\tau = 10^{-12}$,
$\kappa = 0.3$. Percentiles are in reflectance percentage points. Excluded entries are
counted, never replaced by a prescribed value; the constrained retrieval of
Section~\ref{sec:theory}, which does replace them, is reported separately.}
\label{tab:constrained}
\end{table}}{}

\IfFileExists{../figures/emit_quantiles_top_models.png}{%
\begin{figure}[p]
\centering
\includegraphics[width=\textwidth]{emit_quantiles_top_models.png}
\caption{Residual quantile bands across the test set, band by band. The envelopes are
nested percentiles of the residual at each wavelength: $25$--$75$ in red, $5$--$25$ and
$75$--$95$ in blue, $2.5$--$5$ and $95$--$97.5$ in green, $0.5$--$2.5$ and
$97.5$--$99.5$ in grey, with the median as a thin black line. Left column: the direct-flux
residual $Y_2 - \widehat{Y_2}$ in physical units. Centre: the reflectance residual over all
bands. Right: the same residual restricted to the admissible domain with the conditioning
floor, on the same vertical scale, so that the difference between the two columns is the
contribution of the excluded entries.}
\label{fig:bands}
\end{figure}}{}

Two features of Figure~\ref{fig:bands} do not appear in the aggregates. The component
residual of every family is smooth in wavelength and widens where the transmitted flux is
small, so the families differ mainly in the amplitude of a shared shape. The reflectance
residual does not inherit that shape: over all bands its outer envelopes open into vertical
walls at the water-vapour features near 1380 and 1900\,nm and at the long-wave edge, in
every family and at the same wavelengths, because those walls are set by the denominator
and not by the emulator. Under the conditioning floor the walls close, and what remains is
what the component errors predict.
